\section{Related Work}

Our paper extends the long lineage of research on studies of scientists' workflows in two novel ways: 1)~by emphasizing how they envision modern AI helping them, and 2)~by focusing on under-studied physical tasks `beyond the desk.'


\subsection{Studies of Scientists' Workflows}

Scientific workflows and collaborative practices have been long-standing topics of interest in HCI, CSCW, and Science and Technology Studies (STS). For instance, ethnographic studies have examined how scientists conduct their work, from the day-to-day of laboratory life~\cite{latour2013laboratory} to the ``shop work'' and ``shop talk'' through which technical practices, material artifacts, and social processes coordinate to produce scientific knowledge~\cite{lynch1979art}. Viewed through this lens, scientific practice---from the lab to the field---is not only the execution of formal protocols, but also a complex negotiation of material and social conditions. Researchers have documented how tacit knowledge~\cite{collins1974tea}, embodied physical skill~\cite{cetina1999epistemic}, infrastructure~\cite{vertesi2014seamful}, and environmental constraints~\cite{yeh2006butterflynet} shape everything from experimental design to collaboration.

Starting in the 2000s, more contemporary work analyzed scientists' programming and data-centric workflows `at the desk.' Studies have examined how scientists develop and maintain software~\cite{hannah2009how, carver2007software}, rely on computational tools~\cite{huang2025scientists}, manage file dependencies~\cite{gori2020fileweaver}, and organize high performance computing (HPC) scripts, directories, and datasets~\cite{strong2011alamos}.

Although prior work has offered rich accounts of how scientists work to produce scientific knowledge, most of it predates the widespread availability of modern AI, especially its use `beyond the desk' out in lab and field settings. Our interview study advances this lineage of research by \rev{being the first to examine how scientific practitioners engaged in physical lab and field work} imagine modern AI assistance within their day-to-day embodied workflows.

\subsection{Human-Centered AI in Science}

With recent advances in deep learning and generative AI, researchers have explored how these technologies can automate aspects of scientific discovery in tasks like molecular structure prediction, chemical synthesis, and hypothesis generation~\cite{ramos2025review, jumper2021highly, reddy2025towards, baek2024researchagent, rapp2024self}. While pushing the frontiers of what AI can do for science, most efforts center around automating processes of modeling, simulation, and algorithmic optimization~\cite{lu2024ai, jablonka202314} with GPU-enhanced computation.

There is an emerging line of work on robotic arms to automate some physical lab processes like pipetting and diluting microfluids; however, these systems are currently expensive and inflexible, so they are limited more to larger-scale industrial production settings and not seen in most academic labs, where these tasks are done by hand~\cite{robotic_lab_assistants_2024}. \rev{Semi-autonomous UAVs and field robots have considerable potential to support field scientists in tasks such as interacting with remote environments~\cite{outreachrobotics}, terrain imaging~\cite{seifert2019influence}, and environmental monitoring~\cite{koukouvelas2023uav}. Emerging work in robotics for fieldwork has begun to explore opportunities in this direction, highlighting challenges in animal species and individual identification, site access, and data handling~\cite{pringle2025opportunities}. However, existing systems are typically deployed for predefined missions, rather than multi-step or improvisational scientific workflows, and research in this area remains focused on technical and algorithmic capabilities.}

Recognizing the limits of full automation, HCI researchers have focused on Human-AI collaboration, where human scientists drive inquiry and AI assists~\cite{schmidgall2025agent, o2023bioplanner}. These systems support tasks like experimental planning and validation~\cite{schmidgall2025agent, o2023bioplanner}, data processing~\cite{jablonka202314}, and research question ideation (e.g. CoQuest~\cite{liu2024how}, PersonaFlow~\cite{liu2025personaflow}).

Most work here has been on building new system prototypes; to our knowledge, there have been only three prior \emph{studies} of how scientists use AI in their workflows. The first in 2023 presented interviews with professors and tech-industry scientists \rev{(often at more senior levels)} about their desired use cases for AI in science; participants envisioned only AI support for knowledge workflows `at the desk', with potential in science education, data wrangling, literature review, coding, and technical writing~\cite{morris2023scientists}. Two more recent studies at CHI 2025 focused again only on desk-based workflows: one studied how scientists and operations staff use an internal ChatGPT-style chatbot at a U.S. national laboratory (e.g., for writing emails, reports, and manuscripts)~\cite{wagman2025generative}, and the other on how scientists use AI for programming and data analysis~\cite{obrien2025how}.

However, despite the rise of computational techniques, much of modern science still happens in the lab or out in the field, and thus exercise scientists' physical skills and embodied cognitive abilities that go `beyond the desk.' Our work differs from these prior studies by focusing on the perspectives of \rev{scientific practitioners} in improvisational, uncertain, and noisy environments such as wet labs or remote outdoor field sites. \rev{Our study provides a situated, embodied perspective of practitioners ranging from graduate students to experienced full-time lab staff, which complement the findings of prior studies focused on the desk-based computer-centric workflows of scientists~\cite{morris2023scientists,wagman2025generative}.} We extend their findings by situating \rev{scientific practitioners'} perceptions, barriers to adoption, and speculative visions of AI within physical contexts. \rev{Taken together, our findings can be combined with those from prior work on scientists' use of AI at-the-desk to give the field a more holistic view of how lab and field science might be done end-to-end in the future with human-centered AI support.}

\subsection{Human-Centered AI in General Knowledge Work} 

Zooming out farther beyond applications to science, consumer-available AI tools capable of programming, writing, and generating imagery have impacted many types of professional expert stakeholders, especially software engineers~\cite{zheng2025towards, vaithilingam2022expectation, khojah2024beyond, liu2023code}, designers~\cite{subramonyam2025prototyping, zhu2018explainable, shi2023understanding}, educators~\cite{prather2025beyond}, writers~\cite{guo2025from, ippolito2022creative}, artists~\cite{chang2023the, tang2024whats, suh2021ai}, and knowledge workers across areas like law, medicine, and business~\cite{amershi2019guidelines, fok2024marco}. Recent studies have been directed towards understanding how AI could change or improve these domain workflows in terms of productivity, creativity, or collaboration~\cite{vanukuru2025designing, pu2025assistance, suh2024luminate}. Studies have also investigated how professionals use AI as a co-creative or assistive partner, such as surfacing tensions in how writers choose to integrate AI assistance~\cite{guo2025from}.

However, the aforementioned domains and studies typically involve `desk work.' One notable exception~\cite{kernan2023harnessing} examined how LLM-based cognitive assistants can support factory workers in physically-demanding environments, for example, in helping workers resolve mechanical issues with production lines. Similarly, Kernan et al. proposes a set of design guidelines that emphasizes real-time data and domain-specific knowledge integration for AI assistants in manufacturing~\cite{kernan2023harnessing}. While this line of research considers factory-style production labor, our study spotlights AI in the context of the behind-the-scenes physical work that powers science.